\SourceLeaf{296}
\begin{center}
\textbf{NACHLASS UND BRIEFWECHSEL.}

\vspace{1em}

\textbf{[ZUR SPHÄRISCHEN TRIGONOMETRIE.]}

\vspace{0.5em}

\end{center}

\vspace{1em}

\begin{center}[1.]\end{center}

\begin{center}
\textsc{Gauss} an \textsc{Gerling}. Göttingen, 18. Februar 1815.
\end{center}

Recht sehr danke ich Ihnen, lieber \textsc{Gerling}, für die Mittheilungen aus
Ihrer Arbeit über die Trigonometrie. Ihren Beweis für die 4 Gleichungen finde
ich für Ihren Zweck sehr angemessen. Es fehlt ihm weiter nichts zur
Vollständigkeit, als dass noch die Zeichen besonders bewiesen werden müssen,
indem z.\,B. aus
\[
\sin \tfrac12 A^2 \cos \tfrac12 (b+c)^2 = \cos \tfrac12 (B+C)^2 \cos \tfrac12 a^2
\]
eigentlich nichts weiter folgt, als dass
\[
\begin{array}{r@{\ }c@{\ }l}
\sin \tfrac12 A \cos \tfrac12 (b+c) & \text{entweder} & {}=\cos \tfrac12 (B+C) \cos \tfrac12 a \\[4pt]
 & \text{oder} & {}=-\cos \tfrac12 (B+C) \cos \tfrac12 a.
\end{array}
\]

Man kann eine solche Completirung des Beweises aus andern Sätzen ableiten,
wodurch aber das Ganze weitläuftiger wird. Ich habe zuletzt mich immer an
folgende Entwickelungsmethode gehalten:

Ich setze die 5 Gleichungen voraus, da sie in der That alle ohnehin
nothwendig sind:
